\subsection{Coverings}

Let Y be a path-connected, nonempty topological space, and let \((\mathrm{A}_{i})_{i\in\mathrm{I}}\) be a covering of Y by nonempty path-connected subsets, indexed by a totally ordered set I. Let \(X=\bigsqcup_{i\in I}A_{i}\) be the sum space of the family \((\mathrm{A}_{i})\) and let \(f:X\to Y\) be the map induced by the family of canonical injections of each \(A_{i}\) in Y. Suppose that the map f satisfies property (VK). This holds in particular in the following two cases:

\begin{enumerate}
\def\labelenumi{(\roman{enumi})}
\item
  the interiors of the sets \(A_{i}\) , for \(i \in I\) cover Y (cf. IV, p. 402, example 2);
\item
  the space Y is delaceable, as are the spaces \(A_{i}\) , for \(i \in I\) , the family \((\mathrm{A}_{i})_{i \in \mathrm{I}}\) is locally finite, the \(A_{i}\) are closed in Y and their pairwise intersections are locally arcwise connected (cf. IV, p. 399, example 1).
\end{enumerate}

Let J' be the set of triples \((i, i', V)\) , where i and \(i'\) are elements of I and where V is a path-connected component of \(A_i \cap A_{i'}\) . If \(j = (i, i', V) \in J'\) we set \(i_1(j) = i\) , \(i_2(j) = i'\) and \(\overline{j} = (i', i, V)\) . Let J be the subset of J' formed by the triples such that \(i < i'\) . We call the frame of the covering the quiver \(\Gamma\) whose set of vertices is I, whose set of arrows is J, and whose origin and target maps are respectively the maps \(j \mapsto i_1(j)\) and \(j \mapsto i_2(j)\) . We shall identify the graph associated with \(\Gamma\) with the graph \(\widetilde{\Gamma}\) whose set of vertices is I, whose set of arrows is \(J \cup \overline{J}\) , whose origin and target maps are the maps \(j \mapsto i_{1}(j)\) and \(j \mapsto i_{2}(j)\) , and whose involution is the map \(j \mapsto \overline{j}\) .

Denote by \(p_1\) and \(p_2\) the projections of the fibered square \(X \times_{Y} X\) onto X; let \(\textsf{Γ}\) be the frame of the pair \((\varpi(p_1), \varpi(p_2))\) of groupoid morphisms from \(\varpi(\mathrm{X} \times_{\mathrm{Y}} \mathrm{X})\) into \(\varpi(\mathrm{X})\) .

Lemma 1. --- The quiver \(\Gamma\) identifies with a subquiver of \(\textsf{Γ}\) ; the quivers \(\Gamma\) and \(\textsf{Γ}\) are connected.

The path-connected components of X are the \(A_{i}\) , for \(i \in I\) . The path-connected components of \(X \times_{Y} X\) are the \((V \times \{i\}) \times_{Y} (V \times \{i'\})\) , for \((i, i', V) \in J'\) . Consequently, the frame \(\textsf{Γ}\) of the pair \((\varpi(p_1), \varpi(p_2))\) is isomorphic to the quiver whose set of vertices is I, whose set of arrows is \(J'\) , the origin and target maps being respectively the maps \(j \mapsto i_1(j)\) and \(j \mapsto i_2(j)\) . The quiver \(\textsf{Γ}\) is connected (IV, p. 406). Moreover, this description identifies \(\Gamma\) with a subquiver of \(\textsf{Γ}\) . Let us also observe that for every arrow j of \(\textsf{Γ}\) , either \(j \in \mathrm{Fl}(\widetilde{\Gamma})\) , or there exists \(i \in I\) such that \(j = (i, i, A_i)\) . It follows that the map \(\pi_0(\Gamma) \to \pi_0(\textsf{Γ})\) deduced from the injection of \(\Gamma\) into \(\textsf{Γ}\) is bijective, so that \(\textsf{Γ}\) is connected.

For every element \(i\) of I, choose a point \(a(i)\) of \(\mathrm{A}_i\) .

For every element \(j = (i, i', \mathrm{V})\) of \(\mathrm{J}\) , choose a point \(b(j)\) on \(\mathrm{V}\) , a path \(\mathrm{B}_1(j)\) connecting \(b(j)\) to \(a(i)\) on \(\mathrm{A}_i\) , and a path \(\mathrm{B}_2(j)\) connecting \(b(j)\) to \(a(i')\) on \(\mathrm{A}_{i'}\) . Let \(j = (i, i', \mathrm{V})\) be an element of \(\overline{\mathrm{J}}\) ; then \(\overline{j} = (i', i, \mathrm{V})\) belongs to \(\mathrm{J}\) , and one sets \(b(\overline{j}) = b(j)\) , \(\mathrm{B}_1(j) = \mathrm{B}_2(\overline{j})\) and \(\mathrm{B}_2(j) = \mathrm{B}_1(\overline{j})\) . For \(j \in \mathrm{J}' \cup \overline{\mathrm{J}'}\) the paths \(\overline{\mathrm{B}_1(j)}\) and \(\mathrm{B}_2(j)\) on \(\mathrm{Y}\) are juxtaposable. Set

\[
\mathrm{B} (j) = \overline {{\mathrm{B} _ {1} (j)}} * \mathrm{B} _ {2} (j).\tag{3}
\]

This is a path joining \(a(i_1(j))\) to \(a(i_2(j))\) in Y; we have the relation \(\mathrm{B}(\overline{j}) = \overline{\mathrm{B}(j)}\) .

For every \(j = (i, i', \mathrm{V}) \in \mathrm{J}'\) , denote by \(p_{j,1} \colon V \to A_i\) and \(p_{j,2} \colon V \to A_{i'}\) the canonical injections; let us also denote by \(\varphi_j \colon \pi_1(\mathrm{V}, b(j)) \to \pi_1(\mathrm{A}_i, a(i))\) and \(\psi_j \colon \pi_1(\mathrm{V}, b(j)) \to \pi_1(\mathrm{A}_{i'}, a(i'))\) the group homomorphisms defined by

\[
\varphi_ {j} (v) = [ \mathrm{B} _ {1} (j) ] ^ {- 1} (p _ {j, 1}) _ {*} (v) [ \mathrm{B} _ {1} (j) ]
\]

and

\[
\psi_ {j} (v) = [ \mathrm{B} _ {2} (j) ] ^ {- 1} (p _ {j, 2}) _ {*} (v) [ \mathrm{B} _ {2} (j) ],
\]

for \(v \in \pi_1(\mathrm{V}, b(j))\) (cf. IV, p. 407).

Fix an element \(i_0\) of I, as well as a subquiver T in the quiver \(\Gamma\) whose associated graph \(\widetilde{T}\) is a maximal tree of the graph \(\widetilde{\Gamma}\) .

For \(i \in I\) , let \((i_0, j_1, i_1, \ldots, j_n, i)\) be the unique path without backtracking joining \(i_0\) to \(i\) in the tree \(\widetilde{T}\) and let us set

\[
\delta (i) = \left[ \mathrm{B} \left(j _ {1}\right) \right] \left[ \mathrm{B} \left(j _ {2}\right) \right] \dots \left[ \mathrm{B} \left(j _ {n}\right) \right];
\]

it is the class of a path joining \(a(i_0)\) to \(a(i)\) in Y. Denote by \(\alpha_i\) the homomorphism from \(\pi_1(\mathrm{A}_i, a(i))\) into \(\pi_1(\mathrm{Y}, a(i))\) deduced from the canonical injection, and let \(\mu_i: \pi_1(\mathrm{A}_i, a(i)) \to \pi_1(\mathrm{Y}, a(i_0))\) be the group homomorphism defined by

\[
\mu_ {i} (v) = \delta (i) \alpha_ {i} (v) \delta (i) ^ {- 1}.
\]

Finally, let \(\mu\colon\mathrm{F}(\mathrm{J})\to\pi_{1}(\mathrm{Y},a(i_{0}))\) the unique group homomorphism such that one has

\[
\mu (j) = \delta (i _ {1} (j)) [ \mathrm{B} (j) ] \delta (i _ {2} (j)) ^ {- 1}
\]

for all \(j \in J\) . There exists a unique group homomorphism

\[
\mathsf {M} \colon \big (\underset {i \in \mathrm{I}} {\ast} \pi_ {1} (\mathrm{A} _ {i}, a (i)) \big) \ast \mathrm{F} (\mathrm{J}) \to \pi_ {1} (\mathrm{Y}, a (i _ {0}))
\]

which coincides with \(\mu_i\) on \(\pi_1(A_i, a(i))\) , for every \(i \in I\) , and with \(\mu\) on \(F(J)\) .

Let K' be the set of quadruples \((i_{1}, i_{2}, i_{3}, \mathrm{U})\) , where \(i_{1}, i_{2}, i_{3}\) are elements of I and where U is a path-connected component of \(A_{i_{1}} \cap A_{i_{2}} \cap A_{i_{3}}\) . For every element \(k = (i_{1}, i_{2}, i_{3}, \mathrm{U})\) of \(K'\) and every pair \((s, t)\) of elements of \(\{1, 2, 3\}\) we set \(j_{st}(k) = (i_{s}, i_{t}, \mathrm{V})\) , where V is the path-connected component of \(A_{i_{s}} \cap A_{i_{t}}\) that contains U; this is an element of \(J'\) .

Let K be the subset of \(K'\) formed by the quadruples \((i_{1}, i_{2}, i_{3}, \mathrm{U})\) such that \(i_{1} < i_{2} < i_{3}\) . For every element \(k = (i_{1}, i_{2}, i_{3}, \mathrm{U})\) of K, choose a point \(c(k)\) of U, as well as paths \(\mathrm{C}_{12}(k)\) , \(\mathrm{C}_{23}(k)\) and \(\mathrm{C}_{13}(k)\) such that \(\mathrm{C}_{st}(k)\) connects \(c(k)\) to \(b(j_{st}(k))\) in \(A_{i_{s}} \cap A_{i_{t}}\) for \(s, t \in \{1, 2, 3\}\) with s \textless{} t.

Let us then set, for \(k \in \mathbf{K}\) ,

\[
\mathrm{L} _ {1} (k) = \overline {{\mathrm{B} _ {1} (j _ {1 3} (k))}} * \overline {{\mathrm{C} _ {1 3} (k)}} * \mathrm{C} _ {1 2} (k) * \mathrm{B} _ {1} (j _ {1 2} (k)),\tag{4}
\]

\[
\mathrm{L} _ {2} (k) = \overline {{\mathrm{B} _ {2} (j _ {1 2} (k))}} * \overline {{\mathrm{C} _ {1 2} (k)}} * \mathrm{C} _ {2 3} (k) * \mathrm{B} _ {1} (j _ {2 3} (k)),
\]

\[
\mathrm{L} _ {3} (k) = \overline {{\mathrm{B} _ {2} (j _ {2 3} (k))}} * \overline {{\mathrm{C} _ {2 3} (k)}} * \mathrm{C} _ {1 3} (k) * \mathrm{B} _ {2} (j _ {1 3} (k)).
\]

For \(s \in \{1,2,3\}\) , denote by \(\lambda_s(k)\) the class in \(\pi_1(\mathrm{A}_{i_s}, a(i_s))\) of the loop \(\mathrm{L}_s(k)\) .

PROPOSITION 1. --- Let Y be a path-connected topological space and let \((\mathrm{A}_{i})_{i\in\mathrm{I}}\) be a covering of Y by nonempty, path-connected subsets, indexed by a totally ordered set I. Suppose that the canonical map from the sum space of the family \((\mathrm{A}_{i})_{i\in\mathrm{I}}\) into Y satisfies property (VK). Then, the homomorphism M introduced above is surjective and its kernel is the smallest normal subgroup containing the following elements:

\[
\begin{array}{l} \text {(R$_ {1}$)} r _ {1} (j) = j \qquad \text {for } j \text { in } \operatorname{Fl}(\mathsf {T}); \\ \text {(R$_ {2}$)} r _ {2} (j, v) = \varphi_ {j} (v) j \psi_ {j} (v) ^ {- 1} j ^ {- 1} \\ \qquad \qquad \qquad \qquad \qquad \qquad \qquad \qquad \qquad \text {for } j = (i,i^{\prime} ,V)\in \mathsf {J} \text { and } v\in\pi_{1} (V,b(j)); \\ \text {(R$_ {3}$)} r _ {3} (k) = \lambda_ {1} (k) j _ {1 2} (k) \lambda_ {2} (k) j _ {2 3} (k) \lambda_ {3} (k) j _ {1 3} (k) ^ {- 1} \\ \qquad \qquad \qquad \qquad \qquad \qquad \qquad \qquad \text {for } k\in \mathsf {K}. \end{array}
\]

Let X be the topological space that is the sum of the family \((\mathrm{A}_{i})_{i\in\mathrm{I}}\) and let \(f\colon X\to Y\) be the map induced by the family of canonical injections of each \(A_{i}\) in Y. The map f satisfies property (VK) by hypothesis; we shall therefore apply to it th. 1 of no. 1. We thus resume the notation of that section and begin by defining a van Kampen datum for the map f.

The path-connected components of X are the sets \(X_{i} = A_{i} \times \{i\}\) for \(i \in I\) . One thus identifies \(\pi_{0}(X)\) with the set I. For every \(i \in I\) , denote by \(\mathfrak{a}(i)\) the point \((a(i), i)\) of \(X_{i}\) .

Let us set \(Z = X \times_{Y} X\) . The path-connected components of \(Z\) are the sets \(Z_j = (V \times \{i\}) \times_{Y} (V \times \{i'\})\) , where \(j = (i, i', V)\) runs through \(J'\) . One thus identifies \(\pi_0(Z)\) with the set \(J'\) . Let \(J_0\) be the set of elements of \(J\) of the form \((i, i, A_i)\) , so that the family \((J_0, J, \overline{J})\) is a partition of \(J'\) . For \(i \in I\) and \(j = (i, i, A_i) \in J_0\) , set \(b(j) = (a(i), a(i))\) and take for \(\beta_1(j)\) and \(\beta_2(j)\) the class of the constant path at \(a(i)\) . For \(j = (i, i', V) \in J \cup \overline{J}\) , denote by \(b(j)\) the point \(((b(j), i), (b(j), i'))\) of \(Z_j\) and denote by \(\beta_1(j)\) and \(\beta_2(j)\) the classes of paths \(t \mapsto (B_1(j)(t), i)\) and \(t \mapsto (B_2(j)(t), i')\) on \(X\) .

Let us set \(W = X \times_{Y} X \times_{Y} X\) . The path-connected components of W are the sets \(W_k = (U \times \{i_1\}) \times_{Y} (U \times \{i_2\}) \times_{Y} (U \times \{i_3\})\) , where \(k = (i_1, i_2, i_3, U)\) runs through \(K'\) . One thus identifies \(\pi_0(W)\) with the set \(K'\) .

Denote by \(K_{0}\) the set of elements of \(K'\) of the form \(k = (i, i, i, A_{i})\) , for \(i \in I\) . For such an element \(k \in K_{0}\) we set \(\mathsf{c}(k) = (\mathsf{a}(i), \mathsf{a}(i), \mathsf{a}(i))\) and one chooses for \(\gamma_{st}(k)\) the class of the constant path at \((\mathsf{a}(i), \mathsf{a}(i))\) .

Denote by \(K_{1}\) the set of elements of \(K'\) of the form \(k = (i_{1}, i_{2}, i_{3}, V)\) for which the set \(\{i_{1}, i_{2}, i_{3}\}\) has two elements. Let k be an element of \(K'\) of the form \((i, i, i', V)\) , such that \(j = (i, i', V)\) belongs to \(J \cup \overline{J}\) . We then set

\[
\mathsf {c} (k) = ((b (j), i), (b (j), i), (b (j), i ^ {\prime})), \quad \gamma_ {1 2} (k) = (\beta_ {1} (j), \beta_ {1} (j))
\]

and we take for \(\gamma_{13}(k)\) and \(\gamma_{23}(k)\) the class of the constant path at \(\mathsf{b}(j)\) . One defines analogously \(c(k)\) , \(\gamma_{12}(k)\) , \(\gamma_{13}(k)\) , \(\gamma_{23}(k)\) for every element \(k\) of \(\mathrm{K}'\) .

For every element \(k = (i_{1}, i_{2}, i_{3}, \mathrm{U})\) of K, set

\[
\mathsf {c} (k) = ((c (k), i _ {1}), (c (k), i _ {2}), (c (k), i _ {3})) ;
\]

for every pair \((s,t)\) of distinct elements of \(\{1,2,3\}\) , take for \(\gamma_{st}(k)\) the image under the map \(x\mapsto((x,s),(x,t))\) of the class of the path \(\mathrm{C}_{st}(k)\) in \(\mathrm{A}_{i_{s}(k)}\cap\mathrm{A}_{i_{t}(k)}\) .

For every point \(x = (x_1, x_2, x_3)\) of \(X \times_Y X \times_Y X\) and any permutation \(\sigma \in \mathfrak{S}_3\) , set \(\sigma(x) = (x_{\sigma^{-1}(1)}, x_{\sigma^{-1}(2)}, x_{\sigma^{-1}(3)})\) . One thus defines an action of the group \(\mathfrak{S}_3\) on W. For every \(k = (i_1, i_2, i_3, U) \in K'\) and any permutation \(\sigma \in \mathfrak{S}_3\) , set similarly \(\sigma(k) = (i_{\sigma^{-1}(1)}, i_{\sigma^{-1}(2)}, i_{\sigma^{-1}(3)}, U)\) ; we have

\[
\sigma \left(\mathrm{W} _ {k}\right) = \left(\mathrm{U} \times \left\{i _ {\sigma^ {- 1} (1)} \right\}\right) \times_ {\mathrm{Y}} \left(\mathrm{U} \times \left\{i _ {\sigma^ {- 1} (2)} \right\}\right) \times_ {\mathrm{Y}} \left(\mathrm{U} \times \left\{i _ {\sigma^ {- 1} (3)} \right\}\right) = \mathrm{W} _ {\sigma (k)}.
\]

Let \(k = (i_1, i_2, i_3, \mathrm{U})\) be an element of \(\mathrm{K}'\) such that \(i_1, i_2, i_3\) are pairwise distinct. There exists a unique permutation \(\sigma \in \mathfrak{S}_3\) such that \(i_{\sigma^{-1}(1)} < i_{\sigma^{-1}(2)} < i_{\sigma^{-1}(3)}\) , such that \(\sigma(k) = (i_{\sigma^{-1}(1)}, i_{\sigma^{-1}(2)}, i_{\sigma^{-1}(3)}, \mathrm{U})\) belongs to \(\mathrm{K}\) . For \(s \in \{1, 2, 3\}\) , one then sets \(c_s(k) = c_{\sigma(s)}(\sigma(k))\) and \(c(k) = (c_1(k), c_2(k), c_3(k))\) , such that \(c(k) = \sigma^{-1}(c(\sigma(k)))\) . For \((s, t) \in \{(1, 2), (1, 3), (2, 3)\}\) , one defines \(\mathrm{C}_{st}(k) = \mathrm{C}_{\sigma^{-1}(s)\sigma^{-1}(t)}(\sigma(k))\) ; this is a path that connects \(c(k)\) to \(b(j_{\sigma(s)\sigma(t)}(\sigma(k))) = b(j_{st}(k))\) .

Denote by \(g\) the quiver morphism from \(\Gamma\) into \(\textsf{Γ}\) which, to a vertex \(i \in I\) of \(\Gamma\) , associates the vertex \(X_i = A_i \times \{i\}\) of \(\textsf{Γ}\) and, to an arrow \(j = (i, i', V) \in J'\) of \(\Gamma\) , associates the arrow \(Z_j = (V \times \{i\}) \times_Y(V \times \{i'\})\) of \(\textsf{Γ}\) . The map Som( \(g\) ) is bijective; the map Fl( \(g\) ) is injective and the graph \(\Gamma\) is connected (IV, p. 410, lemma 1) and the image under \(g\) of the subquiver T is a subquiver \(\mathsf{T}\) of \(\textsf{Γ}\) whose associated graph is a maximal tree of the graph \(\widetilde{\textsf{Γ}}\) .

The points \(\mathsf{a}(i)\) , for \(i \in I\) , the point \(\mathsf{b}(j)\) , for \(j \in J'\) , the points \(\mathsf{c}(k)\) , for \(k \in K'\) , the path classes \(\beta_{1}(j)\) and \(\beta_{2}(j)\) , for \(j \in J'\) , the path classes \(\gamma_{st}(k)\) , for \(k \in \mathrm{K}'\) , the subquiver \(g(\mathrm{T})\) of \(\textsf{Γ}\) and the element \(i_0\) of I define a van Kampen datum of \(f\) .

Denote by \(\rho\) the unique group homomorphism

\[
\rho \colon \big (\underset {i \in \mathrm{I}} {*} \pi_ {1} (\mathrm{X}, \mathsf {a} (i)) \big) * \mathrm{F} (\mathrm{J} ^ {\prime}) \to \big (\underset {i \in \mathrm{I}} {*} \pi_ {1} (\mathrm{A} _ {i}, a (i)) \big) * \mathrm{F} (\mathrm{J})
\]

which induces the isomorphism of \(\pi_1(\mathrm{X},\mathsf{a}(i))\) onto \(\pi_1(\mathrm{A}_i,a(i))\) deduced from the identification of \(\mathrm{A}_i\times \{i\}\) and \(\mathrm{A}_i\) , for every \(i\in \mathrm{I}\) , and such that one has

\[
\begin{array}{l l} \rho (j) = 1 & \text {for } j \in \mathrm{J} _ {0} \\ \rho (j) = j, \quad \rho (\overline {{j}}) = j ^ {- 1} & \text {for } j \in \mathrm{J}. \end{array}
\]

Let L be the group homomorphism defined in th. 1 of IV, p. 408. For \(j = (i, i, \mathrm{A}_i) \in \mathrm{J}_0\) , we have \(\mathsf{L}(j) = 1 = \mathsf{M} \circ \rho(j)\) . Let \(j = (i, i', \mathrm{V})\) an element of J; we have \(\mathsf{L}(j) = (\mathsf{M} \circ \rho)(j)\) by definition. Finally, if \(j = (i, i', \mathrm{V})\) is an element of \(\overline{\mathbf{J}}\) , \(\overline{j} \in \mathbf{J}\) and one verifies that

\[
\mathsf {L} (j) = \mathsf {L} (\overline {{j}}) ^ {- 1} = \mathsf {M} (\rho (\overline {{j}})) ^ {- 1} = \mathsf {M} (\rho (j)).
\]

Consequently, we have M ∘ ρ = L.

The homomorphism L is surjective (loc. cit.), hence the homomorphism M is also surjective. Since the homomorphism \(\rho\) is surjective, the kernel of M is the smallest normal subgroup of \(\left(\ast\pi_{1}(\mathrm{A}_{i},a(i))\right)*\mathrm{F}(\mathrm{J})\) which contains the images under \(\rho\) of the elements defined by the relations \((\mathrm{R}_{1})\) , \((\mathrm{R}_{2})\) , \((\mathrm{R}_{3})\) of Theorem 1 (IV, p. 408). The proof will be complete once we have verified that these images are, besides the elements defined by the relations \((\mathrm{R}_{1})\) , \((\mathrm{R}_{2})\) , \((\mathrm{R}_{3})\) of prop. 1, the elements conjugate to them, or conjugate to their inverses, as well as the identity element.

Elements R \(_{1}\) --- An arrow of the oriented tree T is of the form Z \(_{j}\) , with \(j = (i, i', V) \in J\) ; its image is the element \(j\) of F(J).

Elements R₂. --- Let \(j = (i, i', \mathrm{V}) \in \mathrm{J}'\) . If \(i = i'\) , we have \(\rho(\mathsf{r}_2(j, v)) = 1\) for all \(v \in \pi_1(\mathrm{A}_i, \mathsf{a}(i))\) . If \(j \in J\) the image of \(\mathsf{r}_2(j, v)\) is the element \(r_2(j, v) = \varphi_j(v) j \psi_j(v)^{-1} j^{-1}\) , for every \(v \in \pi_1(\mathrm{Z}, \mathsf{b}(j))\) . In the remaining case, we have \(\overline{j} \in J\) and the equality

\[
\begin{array}{r l} & {\rho (\mathtt {r} _ {2} (j, v)) = \rho (\varphi_ {j} (v) j \psi_ {j} (v) ^ {- 1} j ^ {- 1})} \\ & {\qquad = [ \mathrm{B} _ {1} (j) ] ^ {- 1} v [ \mathrm{B} _ {1} (j) ] \rho (j) [ \mathrm{B} _ {2} (j) ] ^ {- 1} v ^ {- 1} [ \mathrm{B} _ {2} (j) ] \rho (j) ^ {- 1}} \\ & {\qquad = [ \mathrm{B} _ {2} (\bar {j}) ] ^ {- 1} v [ \mathrm{B} _ {2} (\bar {j}) ] \bar {j} ^ {- 1} [ \mathrm{B} _ {1} (\bar {j}) ] ^ {- 1} v ^ {- 1} [ \mathrm{B} _ {1} (\bar {j}) ] \bar {j}} \end{array}
\]

entails that \(\rho(\mathsf{r}_2(j,v))\) is conjugate to \(\rho(\mathsf{r}_2(\overline{j},v^{-1}))\) .

Elements R₃. — Let \(k = (i_{1}, i_{2}, i_{3}, \mathrm{U})\) be an element of K'.

If \(k \in \mathrm{K}_0\) , \(i_1 = i_2 = i_3\) , \(\lambda_s(k)\) is the class of the trivial path for all \(s \in \{1,2,3\}\) , \(j_{st}(k) \in \mathrm{J}_0\) for every pair \((s,t) \in \{(1,2),(1,3),(2,3)\}\) . Then, \(\rho(\mathsf{r}_3(k))\) is the identity element.

Suppose \(k\in \mathrm{K}_1\) . If \(i_1 = i_2\) then \(j = (i_{1},i_{3},\mathrm{U})\in \mathrm{J}\cup \overline{\mathrm{J}}\) and one has

\[
\mathsf {r} _ {3} (k) = \beta_ {1} (j) ^ {- 1} j _ {1 2} \beta_ {1} (j) j _ {2 3} \beta_ {2} (j) ^ {- 1} \beta_ {2} (j) j _ {1 3} ^ {- 1}
\]

whose image under \(\rho\) is the identity element. The other cases are handled similarly.

Suppose that \(i_{1}, i_{2}, i_{3}\) are pairwise distinct. If \(i_{1} < i_{2} < i_{3}\) , \(k \in \mathbf{K}\) and the image of \(\rho(\mathsf{r}_{3}(k))\) is the element \(r_{3}(k)\) .

Let \(\sigma \in \mathfrak{S}_3\) be the permutation that maps 1 to 2 and 2 to 3. We have

\[
\begin{array}{r l} \lambda_ {1} (\sigma (k)) & = \beta_ {1} (j _ {1 3} (\sigma (k))) ^ {- 1} \cdot p _ {1, *} (\gamma_ {1 3} (\sigma (k))) ^ {- 1} \cdot \\ & \qquad \cdot p _ {1, *} (\gamma_ {1 2} (\sigma (k))) \cdot \beta_ {1} (j _ {1 2} (\sigma (k))) \\ & = \beta_ {1} (j _ {3 2} (k)) ^ {- 1} \cdot p _ {1, *} (\gamma_ {3 2} (k)) ^ {- 1} \cdot p _ {1, *} (\gamma_ {3 1} (k)) \cdot \beta_ {1} (j _ {3 1} (k)) \\ & = \beta_ {2} (j _ {2 3} (k)) ^ {- 1} \cdot p _ {2, *} (\gamma_ {2, 3} (k)) ^ {- 1} \cdot p _ {2, *} (\gamma_ {1 3} (k)) \cdot \beta_ {2} (j _ {1 3} (k)) \\ & = \lambda_ {3} (k). \end{array}
\]

We check similarly that we have \(\lambda_2(\sigma(k)) = \lambda_1(k)\) and \(\lambda_3(\sigma(k)) = \lambda_2(k)\) . Consequently,

\[
\begin{array}{r l} & {\rho (\mathsf {r} _ {3} (\sigma (k))) = \lambda_ {1} (\sigma (k)) \rho (j _ {1 2} (\sigma (k))) \lambda_ {2} (\sigma (k)) \cdot} \\ & {\qquad \cdot \rho (j _ {2 3} (\sigma (k))) \lambda_ {3} (\sigma (k)) \rho (j _ {1 3} (\sigma (k))) ^ {- 1}} \\ & {\qquad = \lambda_ {3} (k) \rho (j _ {3 1} (k)) \lambda_ {1} (k) \rho (j _ {1 2} (k)) \lambda_ {2} (k) \rho (j _ {3 2} (k)) ^ {- 1}} \\ & {\qquad = \lambda_ {3} (k) j _ {1 3} (k) ^ {- 1} \lambda_ {1} (k) j _ {1 2} (k) \lambda_ {2} (k) j _ {2 3} (k),} \end{array}
\]

which proves that \(\rho(\mathsf{r}_3(\sigma(k)))\) is conjugate to

\[
\lambda_ {1} (k) j _ {1 2} (k) \lambda_ {2} (k) j _ {2 3} (k) \lambda_ {3} (k) j _ {1 3} (k) ^ {- 1} = \rho (\mathsf {r} _ {3} (k)).
\]

Let \(\tau \in \mathfrak{S}_3\) be the transposition with support \(\{1,2\}\) . We have

\[
\lambda_ {1} (\tau (k)) = \lambda_ {2} (k) ^ {- 1}, \quad \lambda_ {2} (\tau (k)) = \lambda_ {1} (k) ^ {- 1} \quad \text {and} \quad \lambda_ {3} (\tau (k)) = \lambda_ {3} (k) ^ {- 1}.
\]

The equalities

\[
\begin{array}{c} \rho (\mathsf {r} _ {3} (\tau (k))) = \lambda_ {2} (k) ^ {- 1} \rho (j _ {2 1} (k)) \lambda_ {1} (k) ^ {- 1} \rho (j _ {1 3} (k)) \lambda_ {3} (k) ^ {- 1} \rho (j _ {2 3} (k)) ^ {- 1} \\ = \big (\rho (j _ {2 3} (k)) \lambda_ {3} (k) \rho (j _ {1 3} (k)) ^ {- 1} \lambda_ {1} (k) \rho (j _ {1 2} (k)) \lambda_ {2} (k) \big) ^ {- 1} \end{array}
\]

show that \(\rho(\mathsf{r}_3(\tau(k)))\) is conjugate to the inverse of

\[
\lambda_ {1} (k) \rho (j _ {1 2} (k)) ^ {- 1} \lambda_ {2} (k) \rho (j _ {2 3} (k)) \lambda_ {3} (k) \rho (j _ {1 3} (k)) ^ {- 1} = \rho (\mathsf {r} _ {3} (k)).
\]

Since the group \(\mathfrak{S}_3\) is generated by the permutations \(\tau\) and \(\sigma\) , it follows that, for every \(k \in \mathrm{K}\) and every \(\sigma \in \mathfrak{S}_3\) , \(\rho(\mathsf{r}_3(\sigma(k)))\) is conjugate to \(\rho(\mathsf{r}_3(k))\) or to its inverse.

Proposition 1 is thus proved.

COROLLARY 1. --- Under the hypotheses of Proposition 1, suppose further that, for every \(i \in I\) the image of the homomorphism from \(\pi_1(A_i, a(i))\) into \(\pi_1(Y, a(i))\) deduced from the canonical injection of \(A_i\) in Y is trivial. The homomorphism \(M': F(J) \to \pi_1(Y, a(i_0))\) deduced from M by restriction is then surjective and its kernel is the smallest normal subgroup containing the elements \(j \in Fl(T)\) and the elements \(j_{12}(k) j_{23}(k) j_{13}(k)^{-1}\) for \(k \in K\) .

Let \(\pi\colon (*_{i\in\mathrm{I}}\pi_1(\mathrm{A}_i,a(i)))*\mathrm{F}(\mathrm{J})\to\mathrm{F}(\mathrm{J})\) be the unique homomorphism that induces the trivial homomorphism on each \(\pi_1(\mathrm{A}_i,a(i))\) and the identity on \(\mathrm{F}(\mathrm{J})\) ; it is surjective. Let \(i\in\mathrm{I}\) . The definition of \(\mathsf{M}\) and the hypothesis that the image of the homomorphism from \(\pi_1(\mathrm{A}_i,a(i))\) into \(\pi_1(\mathrm{Y},a(i))\) deduced from the canonical injection is trivial imply that, for every \(v\in\pi_1(\mathrm{A}_i,a(i))\) , \(\mathsf{M}(v)\) is the identity element of \(\pi_1(\mathrm{Y},a(i_0))\) . We therefore have \(\mathsf{M}=\mathsf{M}'\circ\pi\) . Consequently, the homomorphism \(\mathsf{M}'\) is surjective and its kernel is the smallest normal subgroup containing the images under \(\pi\) of the elements \(r_1(j)\) , \(r_2(j,v)\) and \(r_3(k)\) defined in Prop. 1. For \(j\in\mathrm{Fl}(\mathrm{T})\) , we have \(\pi(r_1(j))=j\) . For all \(j\in\mathrm{J}\) and every \(v\in\pi_1(\mathrm{V},b(j))\) , we have \(\pi(r_2(j,v))=e\) . Finally, for every \(k=(i_1,i_2,i_3,\mathrm{U})\in\mathrm{K}\) and every \(s\in\{1,2,3\}\) , \(\lambda_s(k)\) is the class of a loop in \(\mathrm{A}_{i_s}\) ; one therefore has \(\pi(\lambda_s(k))=e\) , so that \(\pi(r_3(k))=j_{12}(k)j_{23}(k)j_{13}(k)^{-1}\) . The corollary follows from this.

COROLLARY 2. --- Under the hypotheses of Proposition 1, suppose further that for every \(i \in I\) , the group \(\pi_1(A_i, a(i))\) is reduced to the identity element and that, for every triple \((i_1, i_2, i_3)\) of pairwise distinct elements of I, the set \(A_{i_1} \cap A_{i_2} \cap A_{i_3}\) is empty. The homomorphism \(\mathsf{M}''\colon \mathrm{F}(\mathrm{J} - \mathrm{Fl}(\mathrm{T})) \to \pi_1(\mathrm{Y}, a(i_0))\) deduced from \(\mathsf{M}\) by restriction is then an isomorphism.

Let \(\pi'\colon\mathrm{F}(\mathrm{J})\to\mathrm{F}(\mathrm{J}-\mathrm{Fl}(\mathrm{T}))\) be the homomorphism which maps \([j]\) onto \([j]\) if \(j\in\mathrm{J}-\mathrm{Fl}(\mathrm{T})\) and which maps \([j]\) to the identity element if \(j\in\mathrm{Fl}(\mathrm{T})\) ; it is surjective and we have \(\mathsf{M}''\circ\pi'= \mathsf{M}'\) , where \(\mathsf{M}'\) is the surjective homomorphism defined in Corollary 1. It follows that \(\mathsf{M}''\) is surjective and its kernel is the smallest normal subgroup of \(\mathrm{F}(\mathrm{J}-\mathrm{Fl}(\mathrm{T}))\) which contains the images under \(\pi'\) of the elements described in loc. cit. Now, by construction, \(\pi'(j) = e\) if \(j \in \mathrm{Fl}(T)\) and the set K is empty, by hypothesis. The homomorphism \(M''\) is therefore an isomorphism.

Examples. --- 1) For the case of a covering formed by two sets, see No. 3.

\begin{enumerate}
\def\labelenumi{\arabic{enumi})}
\setcounter{enumi}{1}
\tightlist
\item
  Let G be a graph (II, p. 155, Definition 1); denote by S the set of vertices of G, by A the set of its oriented edges, and by o and t the origin and terminus maps from A to S; for every oriented edge \(a \in A\) , write \(\overline{a}\) the opposite oriented edge. Endow the sets S and A with the discrete topology; let X be the sum space of the space S and the space I × A, and let ∼ be the finest equivalence relation on X for which \((u, a) \sim (1 - u, \overline{a})\) , \((0, a) \sim o(a)\) and \((1, a) \sim t(a)\) for all \(u \in I\) and every oriented edge \(a \in A\) . The quotient space \(|G| = X / \sim\) is called the geometric realization of the graph G. We denote by p the canonical projection of X onto \(|G|\) .
\end{enumerate}

Let us prove that \(|\mathbf{G}|\) is locally contractible. Let \(s \in S\) . Denote by \(\mathrm{X}_s\) the union of \(\{s\}\) and the subsets \([0,1[ \times \{a\}\) for \(a \in o^{-1}(s)\) and the subsets \(]0,1] \times \{a\}\) for \(a \in t^{-1}(s)\) . Let \(\mathrm{U}_s\) be the image of \(\mathrm{X}_s\) in \(|\mathbf{G}|\) ; it is an open neighborhood of \(p(s)\) in \(|\mathbf{G}|\) since \(\mathrm{X}_s\) is a saturated open neighborhood of \(s\) in \(X\) . Let \(f\) be the map from \(\mathrm{X}_s \times \mathbf{I}\) into \(\mathrm{X}_s\) defined, for \(u, v \in \mathbf{I}\) and \(a \in A\) by the relations

\[
\begin{array}{c} f (s, v) = s \\ f ((u, a), v) = \left\{ \begin{array}{l l} ((1 - v) u, a) & \text {if } 0 \leqslant u <   1 \text {and } o (a) = s, \\ (1 - (1 - v) (1 - u), a) & \text {if } 0 <   u \leqslant 1 \text {and } t (a) = s. \end{array} \right. \end{array}
\]

It is continuous and compatible with the equivalence relation ∼. It therefore defines, by passing to the quotient, a map \(\varphi_{s}\colon U_{s}\times I\to U_{s}\) which is continuous, since I is locally compact (I, p. 19, prop. 10). It is a strong contraction of \(U_{s}\) onto \(p(s)\) (III, p. 237, def. 6).

Let moreover \(x = (\tau, a) \in ]0, 1[ \times \mathrm{A}\) . Denote by \(\mathrm{X}_x = ]0, 1[ \times \{a, \overline{a}\}\) and let \(\mathrm{U}_x\) be its image in \(|\mathbf{G}|\) under \(p\) ; it is a neighborhood of \(p(x)\) in \(|\mathbf{G}|\) , homeomorphic to \(]0, 1[\) . The map from \(\mathrm{X}_x \times \mathbf{I}\) into \(\mathrm{X}_x\) given by \(((u, a), v) \mapsto ((1 - v)u + v\tau, a)\) and \(((u, \overline{a}), t) \mapsto ((1 - v)u + v(1 - \tau), \overline{a})\) is continuous. By passing to the quotient it defines a map \(\varphi_x: \mathrm{U}_x \times \mathbf{I} \to \mathrm{U}_x\) which is continuous (loc. cit.) and is a strong contraction at \(x\) .

Every point of \(|\mathbf{G}|\) is the image of a point \(s \in \mathbf{S}\) or of a point of \(\mathbf{I} \times \mathbf{A}\) of the form \((\tau, a)\) where \(0 < \tau < 1\) . It follows that every point of \(|\mathbf{G}|\) has a contractible neighborhood at this point. In other words, \(|\mathbf{G}|\) is locally contractible and, in particular, simply connected (IV, p. 346, prop. 5).

Choose a total order on S and consider the open covering \((\mathrm{U}_{s})_{s\in\mathrm{S}}\) of \(|G|\) . Let s and \(s'\) be distinct elements of S. The intersection \(U_{s}\cap U_{s'}\) is the union of the \(p([0,1[,a)\) where a ranges over the set of oriented edges of G whose endpoints are s and \(s'\) . Consequently, the nerve of the covering \((\mathrm{U}_{s})_{s\in\mathrm{S}}\) is identified with the quiver G. Moreover, for every \(s\in S\) , \(U_{s}\) is contractible at s, hence \(\pi_{1}(\mathrm{U}_{s},p(s))\) is reduced to the identity element. It then follows from cor. 2 of IV, p. 416 that for every \(s\in S\) , \(\pi_{1}(|\mathrm{G}|,p(s))\) is a free group (Nielsen-Schreier theorem).

COROLLARY 3. --- Under the hypotheses of prop. 1, suppose moreover that the frame of the covering (A \(_{i}\) ) \(_{i\in I}\) is an oriented tree. The homomorphism N: * \(_{i\in I} \pi_{1}(A_{i},a(i)) \to \pi_{1}(Y,a(i_{0}))\) deduced from M by restriction is then surjective; its kernel is the smallest distinguished subgroup of * \(_{i\in I} \pi_{1}(A_{i},a(i))\) which contains the elements \(\varphi_{j}(v)\psi_{j}(v)^{-1}\) , for every \(j = (i_{1},i_{2},V) \in J\) and every \(v \in \pi_{1}(V,b(j))\) .

If the path-connected components of the intersections \(A_{i} \cap A_{i'}\) , for \(i \neq i'\) are moreover simply connected by arcs, the homomorphism N is then an isomorphism.

Under the hypotheses of the corollary, the graph associated with the quiver \(\Gamma\) is a tree, whence \(T = \Gamma\) . It follows that the image of the group \(F(J)\) under the homomorphism \(M\) is reduced to the identity element.

Let

\[
\rho \colon \big (\underset {i \in \mathrm{I}} {\ast} \pi_ {1} (\mathrm{A} _ {i}, a (i)) \big) \ast \mathrm{F} (\mathrm{J}) \to \underset {i \in \mathrm{I}} {\ast} \pi_ {1} (\mathrm{A} _ {i}, a (i))
\]

the unique group homomorphism that induces the identity homomorphism on \(\pi_1(A_i, a(i))\) and whose kernel contains F(J). We have M = N \(\circ\) \(\rho\) . Consequently, the homomorphism N is surjective and its kernel is the smallest distinguished subgroup of \(*_{i \in I} \pi_1(A_i, a(i))\) which contains the images under \(\rho\) of the elements defined by the relations (R \(_1\) ), (R \(_2\) ) and (R \(_3\) ) of prop. 1. We have \(\rho(j) = 1\) for all \(j \in \text{Fl}(\Gamma)\) . Since the frame of the covering (A \(_i\) ) \(_{i \in I}\) is a tree, we have A \(_{i_1} \cap\) A \(_{i_2} \cap\) A \(_{i_3} = \varnothing\) for every triple \(i_1, i_2, i_3\) ) of distinct elements of I. It follows that the kernel of N is the smallest distinguished subgroup containing the elements \(\varphi_j(v) \psi_j(v)^{-1}\) for all \(j = (i_1, i_2, V) \in J\) and every \(v \in \pi_1(V, b(j))\) .

Example 3 (Private plane of n points). --- The fundamental group of \(R^{2}-\{0\}\) is isomorphic to Z and the class of the path \(t \mapsto e^{2\pi it}\) it is a generator (IV, p. 347, corollary). More generally, let n be a natural number and let \(A = \{z_{1}, \ldots, z_{n}\}\) a set of n points of \(R^{2}\) Let Y be the space \(R^{2}-A\) ; we shall prove that the fundamental group of Y is isomorphic to the free group \(F_{n}\) on n generators. For every i, set \(z_{i} = (u_{i}, v_{i})\) . After possibly replacing \(z_{i}\) under \(f(z_{i})\) , where \(f: R^{2} \to R^{2}\) is a homeomorphism of the form \((u, v) \mapsto (u + \alpha v, v)\) , we may assume that the abscissas of the \(z_{i}\) are pairwise distinct. It is also not restrictive to assume that we have \(u_{1} < \cdots < u_{n}\) .

Let us set \(V_{1} = ]-\infty, u_{2}[ \times R\), \(V_{i} = ]u_{i-1}, u_{i+1}[ \times R\) for \(2 \leqslant i \leqslant n-1\) , \(V_{n} = ]u_{n-1}, +\infty[ \times R\) . For \(1 \leqslant i \leqslant n\) the set \(U_{i} = V_{i} - \{z_{i}\}\) is open in the plane and homeomorphic to \(R^{2} - \{0\}\) . The family \((U_{i})_{1 \leqslant i \leqslant n}\) is an open cover of the space \(Y = R^{2} - A\) . The intersection \(U_{i} \cap U_{j}\) is empty for \(|i - j| \geqslant 2\) , homeomorphic to \(R^{2}\) for \(|i - j| = 1\) . By Corollary 3, the fundamental group of Y is isomorphic to the free group \(F_{n}\) .

Let \(a\) a point of Y. For every integer \(i \in \{1, \ldots, n\}\) , let \(r_i\) be a strictly positive real number strictly less than the distances from \(z_i\) to the points \(z_j\) , for \(j \neq i\) . Denote by \(v_i\) the class of the loop \(t \mapsto z_i + r_i e^{2\pi it}\) at the point \(z_i + r_i\) of Y. Let \(\theta_i\) be the class of a path \(\gamma_i\) connecting the point \(a\) at the point \(z_i + r_i\) in Y. If the paths \(\gamma_i\) are injective and that their images meet only at \(a\) , one can prove that the unique homomorphism from the free group F( \(t_1, \ldots, t_n\) ) in \(\pi_1(Y, a)\) such that \(\varphi(t_i) = \theta_i v_i \theta_i^{-1}\) for all \(i \in \{1, \ldots, n\}\) is a group isomorphism.

An analogous argument allows one to show that for any closed discrete subset A of the plane, the fundamental group of \(R^{2}-A\) is isomorphic to F(A) (IV, p. 463, exerc. 1).

Corollary 4. --- Under the hypotheses of Proposition 1, suppose that there exists a nonempty path-connected subset A of Y such that the intersection \(\mathrm{A}_i \cap \mathrm{A}_{i'}\) be equal to A for every pair \((i, i')\) of distinct elements of I. Let a be a point of A. There exists a unique homomorphism \(\varphi\) of the sum of the family of groups \((\pi_1(\mathrm{A}_i, a))_{i \in \mathrm{I}}\) amalgamated by \(\pi_1(\mathrm{A}, a)\) on \(\pi_1(\mathrm{Y}, a)\) which coincides with the homomorphism deduced from the canonical injection of \(\mathrm{A}_i\) in Y, for every \(i \in \mathrm{I}\) . The homomorphism \(\varphi\) is an isomorphism.

In particular, if the group \(\pi_{1}(\mathrm{A},a)\) is reduced to the identity element, the canonical homomorphism of the free product of the family of groups \((\pi_1(\mathrm{A}_i,a))_{i\in \mathrm{I}}\) on \(\pi_1(\mathrm{Y},a)\) is an isomorphism.

For \(i \in \mathrm{I}\) , denote by \(g_i \colon \pi_1(\mathrm{A}, a) \to \pi_1(\mathrm{A}_i, a)\) and \(f_i \colon \pi_1(\mathrm{A}_i, a) \to \pi_1(\mathrm{Y}, a)\) the canonical homomorphisms deduced from the inclusions of \(\mathrm{A}\) on \(\mathrm{A}_i\) and of \(\mathrm{A}_i\) on \(\mathrm{Y}\) . Denote by \(*_{\mathrm{A}} \pi_1(\mathrm{A}_i, a)\) the sum of the groups \(\pi_1(\mathrm{A}_i, a)\) amalgamated by \(\pi_1(\mathrm{A}, a)\) Let also \(p\) the unique homomorphism from \(*_{i \in \mathrm{I}} \pi_1(\mathrm{A}_i, a)\) on \(*_{\mathrm{A}} \pi_1(\mathrm{A}_i, a)\) which induces the identity on \(\pi_1(\mathrm{A}_i, a)\) (A, I, p. 80, prop. 4). The homomorphism \(p\) is surjective, and it follows from the definition of the monoid \(*_{\mathrm{A}} \pi_1(\mathrm{A}_i, a)\) that its kernel is the smallest normal subgroup of \(*_{i \in \mathrm{I}} \pi_1(\mathrm{A}_i, a)\) which contains the elements \(g_i(v)g_{i'}(v)^{-1}\) , for \(i, i' \in \mathrm{I}\) and \(v \in \pi_1(\mathrm{A}, a)\) .

Homomorphisms \(f_i \circ g_i \colon \pi_1(A, a) \to \pi_1(Y, a)\) , for \(i \in I\) , are equal. It therefore follows from the universal property of amalgamated sums of monoids (A, I, p. 80, prop. 4) that there exists a unique group homomorphism \(\varphi\) (resp. \(f\) ) of \(*_{A} \pi_1(A_i, a)\) (resp. of \(*_{i \in I} \pi_1(A_i, a)\) ) in \(\pi_1(Y, a)\) which induces the homomorphism \(f_i\) onto \(\pi_1(A_i, a)\) . We have \(f = \varphi \circ p\) .

For \(i \in I\) , denote by \(u_i\) the canonical injection of A into \(A_i\) and \(v_i\) the canonical injection of \(A_i\) in Y. We also denote \(w\) the canonical injection of A into Y. For every \(i \in I\) , we have \(v_i \circ u_i = w\) , therefore \(\pi_1(v_i, a) \circ \pi_1(u_i, a) = \pi_1(w, a)\) It then follows from the universal property of amalgamated sums of monoids (A, I, p. 80, prop. 4) that there exists a unique group homomorphism \(\varphi\) of \(*_{i} \pi_1(A_i, a)\) on \(\pi_1(Y, a)\) which induces the homomorphism \(\pi_1(v_i, a)\) onto \(\pi_1(A_i, a)\) . It is a matter of proving that \(\varphi\) is an isomorphism.

The set J is identified with the set of pairs \((i, i')\) of elements of I such that \(i < i'\) . The set K is identified with the set of triples \((i_{1}, i_{2}, i_{3})\) of elements of I such that \(i_{1} < i_{2} < i_{3}\) Let us choose all base points \(a(i)\) , \(b(j)\) and \(c(k)\) equal to a and all paths \(\mathrm{B}(j)\) , \(\mathrm{C}_{st}(k)\) equal to the constant path with image a. We also fix a point \(i_{0} \in I\) .

For every pair \((i, i')\) of distinct points of I, the skeleton \(\Gamma\) of the covering \((\mathrm{A}_{i})\) has exactly one arrow with endpoints i and \(i'\) The arrows of \(\Gamma\) one of whose endpoints is equal to \(i_{0}\) are the arrows of a maximal oriented subtree T.

For every element \(k \in \mathrm{K}\) , we have \(r_3(k) = j_{12}(k)j_{23}(k)j_{13}(k)^{-1}\) ; let R be the distinguished subgroup of \(\mathrm{F(J)}\) generated by the elements \(j \in \mathrm{Fl(T)}\) and the elements \(r_3(k)\) , \(k \in \mathrm{K}\) . Let us prove that \(\mathrm{R} = \mathrm{F(J)}\) . It suffices to show that every element \(j = (i_1, i_2)\) of J belongs to R. This is true, by hypothesis, if \(i_1 = i_0\) or \(i_2 = i_0\) . Suppose \(i_0 < i_1\) and let us set \(k = (i_0, i_1, i_2)\) . It is an element of K such that \(j = j_{23}(k) = j\) Moreover, \(j_{12}(k)\) and \(j_{13}(k)\) belong to Fl(T). It follows that \(j\) belongs to R. The cases where \(i_1 < i_0 < i_2\) or \(i_2 < i_0\) are treated in an analogous manner.

It then follows from prop. 1 (IV, p. 412) that the homomorphism \(f\) is surjective and its kernel is the smallest normal subgroup containing the relators \(r_2(j,v) = g_i(v)g_{i'}(v)^{-1}\) , for \(j = (i,i',\mathrm{A}) \in \mathrm{J}\) and \(v \in \pi_1(\mathrm{A},a)\) In other words, \(\operatorname{Ker}(f) = \operatorname{Ker}(p)\) This implies that \(\varphi\) is an isomorphism, as was to be shown.

If \(\pi_{1}(A,a)\) is reduced to the identity element, p is an isomorphism, whence the second assertion.

Example 4. — Let \(((\mathrm{X}_{i}, x_{i}))_{i \in \mathrm{I}}\) be a family of pointed topological spaces. We call the wedge sum of the family \(((\mathrm{X}_{i}, x_{i}))_{i \in \mathrm{I}}\) , and denote by \(\bigvee_{i \in \mathrm{I}} (\mathrm{X}_{i}, x_{i})\) the quotient topological space of the disjoint union of the family \((\mathrm{X}_{i})_{i \in \mathrm{I}}\) by the equivalence relation which identifies all the points \((x_{i}, i)\) with one another, for \(i \in \mathrm{I}\) . Denote this topological space by X and the common image of the \(x_{i}\) by x. Suppose that, for every \(i \in \mathrm{I}\) , the point \(x_{i}\) is closed in \(\mathrm{X}_{i}\) and that the spaces \(\mathrm{X}_{i}\) are delaceable. If I is finite, Corollary 4 implies that the canonical homomorphism \(* \pi_{1}(\mathrm{X}_{i}, x_{i}) \to \pi_{1}(\mathrm{X}, x)\) is an isomorphism.

Remark 1 of IV, p. 429 and Exercise 3, IV, p. 463 give less restrictive conditions under which this homomorphism is an isomorphism. See nevertheless Exercise 4, IV, p. 464.

